% =====================================================================
%  后记（あとがき，原书 p.176–187 / PDF 185–196）
% =====================================================================
\chapter*{后记}
\addcontentsline{toc}{chapter}{后记}
\markboth{后记}{后记}

% [原文 p.176]
如前所述，按最初的计划，除上述内容之外，本还打算从与辛几何相关的话题中，选取若干近年来因与数学物理及超弦理论的关联而受人注目的内容加以介绍，但因篇幅所限未能实现。在这篇后记中，想就这些话题稍作评述。

但在那之前，先把与正文所述内容相关的参考文献列出。第 2 章沿着一条路径追踪了从分析力学通向辛几何的脉络，但这假定读者事先已具备分析力学最低限度的知识。不懂分析力学的人若想在这一章里临时抱佛脚地学分析力学，那是鲁莽的。所幸分析力学从浅显的入门书到发展性的专著，各种教科书与参考书应有尽有，读者不妨去书店或图书馆，尽情寻找合乎自己口味的一本。公认评价较高的教科书，可以举出 Landau–Lifshitz《理论物理学教程》中的《力学》。原著是 [85]，日译本由东京图书出版。此外，从数学立场写的教科书有 Arnold《经典力学的数学方法》[9]，日译本由岩波书店出版。
% [Tutor Note: 背景。原书在此只提「邦訳」（日译本）；中译本读者更常用的是中译本：Landau–Lifshitz《力学》（理论物理学教程第一卷，高等教育出版社有译本），Arnold《经典力学的数学方法》（科学出版社／高等教育出版社有译本）。读者若找不到日译本，可径用中译本对照，公式编号一致。]

又假定读者已经学过流形的初步。流形虽是自然的概念，却不能说是所有理工科学生都熟悉的东西。第一道关口是定义，这归根结底只能硬记。第二道（大概也是最大的）关口是切空间，一旦弄懂（或者至少自以为弄懂），流形入门课程的关口就算过了。切空间之后的微分形式不过是完全反对称张量的速记法。入门课程的压轴是 Stokes 定理，它不过是微积分基本定理的改写，却把 Green 定理、Cauchy 积分定理、Stokes 旋转定理、Gauss 散度定理等作为特例包含在内。初次学习微分形式的人，到此定会惊叹微分形式提供了何等便利的记法。又，Stokes 定理主张边缘算子与外微分算子互为伴随算子，由此自然联系到 de Rham 定理。
% [Tutor Note: 难点。Stokes 定理 $\int_M \rd\omega = \int_{\partial M}\omega$ 的「伴随」说法，指的是 $(d\omega, \eta) = (\omega, \delta\eta)$ 这一形式伴随关系：外微分 d 与余微分 $\delta$ 在 $L^2$ 内积下互为伴随，而 $\delta$ 正是边缘算子 $\partial$ 在链复形层面的对偶。原书此处说的是「边界作用素」与「外微分作用素」互为伴随作用素，宜理解为：链层面 $\partial$ 与上链层面 $d$ 的对偶配对，以及 Hodge 星号把两者联系起来。]

关于流形基础事项、可供通读的文献，可以举出
\begin{quote}
\noindent [森田 05] 森田茂之，《微分形式的几何学》，岩波书店，2005 年
\end{quote}
这本书在页数与内容上的平衡非常出色。另一方面，若希望关于初等定义有更恳切详尽的说明，则宜读
\begin{quote}
\noindent [松本 88] 松本幸夫，《流形的基础》，东京大学出版会，1988 年
\end{quote}
% [Tutor Note: 背景。森田茂之《微分形式の幾何学》与松本幸夫《多様体の基礎》都是日本流形论的标准入门书，前者以 de Rham 理论见长，后者以定义细致、图多著称。]

关于流形论初步的英文文献极多，评价好的教科书可以举出 Tu [130]。附带一提，同一位作者与 Bott 的合著 [26] 作为用微分形式讲代数拓扑的教科书很有名。从初步讲到中级程度内容的流形教科书，Warner [135] 是常备的一种。这本书的特色，当数用层来证明 de Rham 定理，以及带证明地讲解 Hodge 理论。不过，层的理论从其起源来自多复变函数论这一点来考虑，也宜与复几何一并学习。层极为有用，却是抽象度很高的概念，初学时的心理障碍未必低。标准内容无论怎么学都差别不大，这里举出 Griffith–Harris [54]。日文文献方面，例如
\begin{quote}
\noindent [小林 05] 小林昭七，《复素几何》，岩波书店，2005 年
\end{quote}
也简洁地讲述了层及其上同调。层与其说是实体，不如说是语言，因此不建议在没有具体例的情况下学一般论。
% [Tutor Note: 校勘。原书正文作「複素幾何」（复（素）几何）；中译名「复素几何」中的「素」为日文「複素＝complex」的直译，中文习惯作「复几何」。此处照录原书名，正文行文一律用「复几何」。]

带证明地介绍 Hodge 理论的书，另外还有
\begin{quote}
\noindent [今野 13] 今野宏，《微分几何》，东京大学出版会，2013 年
\end{quote}
本书几乎没有述及的椭圆型估计，也可以在这本书里读到。这本书相对于页数而言内容极为丰富，读过一遍便可具备一整套常识。
% [原文 p.177]
第 3 章到第 5 章选取了辛几何的标准话题。谈到辛几何方面可读的日文文献，无论如何都要举出
\begin{quote}
\noindent [深谷 08] 深谷贤治，《辛几何学》，岩波书店，2008 年
\end{quote}
这本书写于 1999 年，后半部处理了当时最前沿的话题，是一本有雄心的书。另外，[今野 13] 也用不多的页数简洁地介绍了辛几何。英文文献则可举出 McDuff–Salamon [90]。更发展性的话题由同两位作者的 [91] 作了讲解。此外，Cannas da Silva [27] 也是一本易读的入门书。

第 6 章所处理的完全可积系统，与其说是辛几何的一个分支，不如说它自身构成了一个有广度与深度的独立领域。完全可积系统大致分为经典可积系统与量子可积系统，二者又各自分为有限自由度的与无限自由度的。与辛几何关系最直接的是有限自由度的经典可积系统，但以 KdV 方程为代表的无限自由度可积系统，以及以 XXX 模型为代表的量子可积系统，也与辛几何有种种关联。关于完全可积系统的文献浩如烟海，这里根本介绍不完。撰写第 6 章时参考了 Reyman–Semenov–Tian–Shansky [110]、Audin [13]，以及
\begin{quote}
\noindent [高崎 01] 高崎金久，《可积系统的世界》，共立出版，2001 年
\end{quote}
Audin [13] 有共立出版出的日译本。关于 KdV 方程，可以举出
\begin{quote}
\noindent [三伊神 07] 三轮哲二、伊达悦朗、神保道夫，《孤子的数理》，岩波书店，2007 年
\end{quote}
此外，关于 Lie 群与 Lie 环参考了
\begin{quote}
\noindent [小大 05] 小林俊行、大岛利雄，《Lie 群与表现论》，岩波书店，2005 年
\end{quote}
关于 Bruhat 分解则参考了
\begin{quote}
\noindent [寺原 06] 寺田至、原田耕一郎，《群论》，岩波书店，2006 年
\end{quote}
% [Tutor Note: 校勘。原书此处作「第 5 章を書く際には」；但所列 [110]、[13]、[高崎 01] 均为完全可积系统（第 6 章）的文献，且上下文正讨论第 6 章，故疑为「第 6 章」之误。]

第 7 章所介绍的 KAM 理论，其创始人之一 Arnold 的解说见 [9]。要更深入地学习动力系统，以
\begin{quote}
\noindent [久矢 06] 久保泉、矢野公一，《力学系》，岩波书店，2006 年
\end{quote}
为良好的出发点。Birkhoff–Khinchin 遍历定理的证明，例如可在
\begin{quote}
\noindent [小谷 05] 小谷真一，《测度与概率》，岩波书店，2005 年
\end{quote}
读到。

% [原文 p.178]
第 8 章所处理的辛环面流形，[今野 13] 中也有解说。英文文献可以举出 Audin [15]。另外，Cannas da Silva [28] 中有解说文章。从代数几何观点处理环面簇的日文书，则以
\begin{quote}
\noindent [小田 85] 小田忠雄，《凸体与代数几何》，纪伊国屋书店，1985 年
\end{quote}
为经典。较近的入门书则有
\begin{quote}
\noindent [石田 00] 石田正典，《环面簇入门》，朝仓书店，2000 年
\end{quote}
英文方面 Fulton [48] 读起来较容易。更网罗性的文献，最近 Cox–Little–Schenck [32] 已经出版。

关于联络、曲率等构成规范理论基础的概念，[小林 89] 有详尽的解说。关于以 Uhlenbeck 定理为代表的规范理论解析侧面，请参照 Donaldson–Kronheimer [37]、[深谷 95]、[桥本 15]。

那么，下面就说说正文中未能涉及的发展性话题。第 2 章所涉话题之中，只有三体问题历经三百余年，至今仍被活跃地研究。作为近年论文的例子，可以举出 Šuvakov–Dmitrašinović [123] 与 Albers–Frauenfelder–van Koert–Paternain [4]。

第 3 章的定理 3.2.4 讲述了 Weinstein 定理：对任意辛流形 $M$ 的任意 Lagrange 子流形 $L$，$L$ 在 $M$ 中的邻域与 $T^*L$ 中零截面的邻域辛同胚。对于这样的 $L$ 的邻域，Arnold [8] 猜想：落入其中的任意紧完全 Lagrange 子流形，都可以通过 $M$ 的 Hamilton 同胚移到 $L$ 上；这被称为近邻 Lagrange 子流形猜想（nearby Lagrangian conjecture）。关于这一猜想的近期进展，例如见 [1], [46], [47], [103], [104]。

第 4 章所处理的 Liouville 区域，是辛流形的基本构成要素。辛形式为恰当的这一条件极强，因此与紧辛流形相比，Liouville 区域远为容易处理。Liouville 向量场外向这一条件，被用来保证伪全纯曲线（它是作用泛函的梯度流）不逃出边界。关于 Liouville 区域的基本问题，是分类其中所含的完全 Lagrange 子流形。研究 Liouville 区域的工具，有辛上同调、缠绕 Fukaya 范畴（wrapped Fukaya category）等，关于这些请参照 [111], [121] 及 [2]。Stein 区域是 Liouville 区域中性质尤佳者，是多复变函数论中的基本对象。关于多复变函数论的教科书为数众多，这里举出
\begin{quote}
\noindent [辻 12] 辻元，《复流形论讲义》，サイエンス社（临时别册·数理科学 SGC-94），2012 年
\end{quote}
作为 Liouville 区域边界出现的接触流形，例如

% [原文 p.179]
\begin{quote}
\noindent [三松 01] 三松佳彦，《3 次元接触构造的拓扑》，日本数学会，2001 年
\end{quote}
可以从日本数学会的网页下载。
% [Tutor Note: 校勘。原书此句为「[三松 01] が日本数学会のホームページからダウンロードできる」，指该文（数理解析研究所講究録别册? 实为日本数学会「数学通信」／「季刊·数学」所载解说）可从学会网页下载 PDF，而非指书。]

第 5 章所介绍的动量映射与辛商，其与几何不变式论（geometric invariant theory）的关系很重要。关于几何不变式论的基本文献是 [100]。可作为日文阅读的个性文献有
\begin{quote}
\noindent [向井 08] 向井茂，《模空间理论 1,2》，岩波书店，2008 年
\end{quote}
\begin{quote}
\noindent [中岛 08] 中岛启，《非线性问题与复几何学》，岩波书店，2008 年
\end{quote}
其第 6 章解说了几何不变式论与动量映射的关系。

第 8 章最后提到的三维普通二重点有锥流形（conifold）这一别名，被弦论家反复提起。从这个奇点造出光滑空间，至少有三条途径：作奇点解消（resolution of singularity）、作复结构的形变（deformation of complex structure），以及取链环（link）。

% [原文 p.180]
有奇点的流形 $X$ 的奇点解消，指的是从光滑流形 $Y$ 到 $X$ 的固有映射 $\varphi:Y\to X$，使得在去掉奇点之后的集合上是同构。奇点解消 $\varphi:Y\to X$ 若满足 $\varphi^*\omega_X=\omega_Y$，则称为 crepant。这是 $Y$ 成为 Calabi–Yau 流形的条件。如第 8 章所见，三维普通二重点有 crepant 奇点解消，此时的 $Y$ 由 $\PP^1$ 上两个同次直线丛 $O_{\PP^1}(-1)\oplus O_{\PP^1}(-1)$ 的直和的全空间给出。
% [Tutor Note: 校勘。原书此处作「同次直線束」，扫描件字迹与 OCR 均不清（OCR 读作「同音」）；「同次直線束」即齐性直线丛（homogeneous line bundle），$\PP^1$ 上的 $\OO(-1)$ 正是齐性直线丛，故译作「同次直线丛」；第 8 章等处译作「重言线丛」，二者同指。]

复结构的形变通过对定义多项式的变形得到。三维普通二重点由 Morse 型多项式 (8.4.91) 定义，引入正数 $\varepsilon$ 后可将其变形为
\begin{equation*}
Y_\varepsilon:=\bigl\{(w_1,w_2,w_3,w_4)\in\C^4\ \big|\ w_1^2+w_2^2+w_3^2+w_4^2=\varepsilon\bigr\}
\tag{1}
\end{equation*}
对 $i=1,2,3,4$ 令 $w_i=a_i+\sqrt{-1}\,b_i$，把 $Y_\varepsilon$ 的定义方程按实部与虚部分开，得
\begin{align*}
a_1^2+a_2^2+a_3^2+a_4^2-b_1^2-b_2^2-b_3^2-b_4^2&=\varepsilon, \tag{2}\\
a_1b_1+a_2b_2+a_3b_3+a_4b_4&=0. \tag{3}
\end{align*}
这可以用
\begin{equation*}
\mathbf{a}=(a_1,a_2,a_3,a_4),\qquad \mathbf{b}=(b_1,b_2,b_3,b_4)
\tag{4}
\end{equation*}
改写成
\begin{equation*}
|\mathbf{a}|^2=|\mathbf{b}|^2+\varepsilon,\qquad \mathbf{a}\cdot\mathbf{b}=0.
\tag{5}
\end{equation*}
特别地 $|\mathbf{a}|\neq0$，于是可定义映射
\begin{equation*}
\pi:Y_\varepsilon\longrightarrow S^3,\qquad (\mathbf{a},\mathbf{b})\longmapsto \frac{\mathbf{a}}{|\mathbf{a}|}.
\tag{6}
\end{equation*}
对任意 $x\in S^3$，此映射的纤维
\begin{equation*}
\pi^{-1}(x)\cong\{\mathbf{b}\in\R^4\mid x\cdot\mathbf{b}=0\}
\tag{7}
\end{equation*}
与 $S^3$ 在 $x\in S^3$ 处的切空间 $T_xS^3$ 自然等同。因此 $Y_\varepsilon$ 同构于 $S^3$ 的切丛的全空间 $TS^3$。切丛平凡的多流形称为可平行化（parallelizable）的。$S^3$ 是 Lie 群，故可平行化，从而 $TS^3$ 与 $S^3\times\R^3$ 同构。在 $Y_\varepsilon$ 上取由 $\C^4$ 的 Euclid 度量诱导的 Kähler 结构，则 $Y_\varepsilon$ 的实部 $Y_\varepsilon\cap\R^4$ 成为与 $S^3$ 同胚的 Lagrange 子流形。因此由 Weinstein 的 Lagrange 近邻定理，$Y_\varepsilon$ 中 $Y_\varepsilon\cap\R^4$ 的邻域与 $T^*S^3$ 辛同胚。用更精细的论证，还可以证明 $Y_\varepsilon$ 与 $T^*S^3$ 辛同胚。

% [Tutor Note: 背景。(1)–(7) 是 conifold 的「形变」描述：$Y_\varepsilon$ 是以 $S^3$ 为底的 $\R^3$ 丛，即 $TS^3$，它与 $T^*S^3$ 辛同胚——这正是「$T^*S^3$ 的零截面在 $S^3$ 的 conifold 形变中」的确切含义。$\varepsilon\to0$ 时 $Y_\varepsilon$ 退化到奇点 $w_1^2+\cdots+w_4^2=0$。]

对原点有孤立奇点的超曲面奇点
\begin{equation*}
Y=\{z\in\C^n\mid f(z)=0\}
\tag{8}
\end{equation*}
而言，$Y$ 与
\begin{equation*}
S_\varepsilon^{2n-1}=\{z\in\C^n\mid |z|=\varepsilon\}
\tag{9}
\end{equation*}
的公共部分称为 $Y$ 的链环（link）。三维普通二重点的链环由
\begin{equation*}
\{\mathbf{a}+\sqrt{-1}\mathbf{b}\in\C^4\mid |\mathbf{a}|=|\mathbf{b}|=\varepsilon/2,\ \mathbf{a}\cdot\mathbf{b}=0\}
\end{equation*}
给出。这与 $S^3$ 的单位切丛（unit tangent bundle）同构，而 $S^3$ 可平行化，故与 $S^3\times S^2$ 微分同胚。

% [原文 p.181]
't Hooft [124] 指出：在具有耦合常数 $g$ 的 $U(N)$ 规范理论中，固定 $g^2N$ 而令 $N$ 变大的极限，看上去像以 $1/N$ 为耦合常数的弦理论。Maldacena（[86]）接过 't Hooft 的想法，猜想四维 $\mathcal{N}=4$ 超 Yang–Mills 理论的巨大 $N$ 极限，与 $AdS_5\times S^5$ 上的 IIB 型超弦理论对偶。这里 $AdS_5$ 是五维反 de Sitter 空间（anti-de Sitter space）。这一猜想被称为 AdS/CFT 对应（AdS/CFT correspondence），引起了巨大关注。关于 AdS/CFT 对应的解说，例如有 Aharony–Gubser–Maldacena–Ooguri–Oz [3]。

在 AdS/CFT 对应中，作为与五维反 de Sitter 空间取直积的 Riemann 流形，可以不用五维球面 $S^5$，而考虑称为佐佐木–Einstein 流形的更广一类流形。这里，Riemann 流形 $(S,g)$ 是佐佐木流形，指的是其 Riemann 锥 $(\R_{>0}\times S,\ dr^2+r^2g)$ 为 Kähler 流形。满足 Einstein 方程的佐佐木流形称为佐佐木–Einstein 流形。这等价于其 Riemann 锥为 Ricci 平坦 Kähler 流形。
% [Tutor Note: 背景。佐佐木流形（Sasaki manifold）的 Riemann 锥为 Kähler 这一等价，是 Sasaki–Einstein 几何的出发点；$S^5$ 的锥就是 $\C^3\setminus\{0\}$，故 $S^5$ 与 $S^2\times S^3$ 的圆锥都对应 Kähler 锥，这正是 AdS/CFT 里把 $S^5$ 换成 Sasaki–Einstein 流形的几何依据。]

三维普通二重点的链环可以等同于 $(SU(2)\times SU(2))/U(1)$，并带有自然的佐佐木–Einstein 度量。这里分母的 $U(1)$ 是把 $SU(2)$ 的极大环面按对角方式嵌入 $SU(2)\times SU(2)$ 所得者。与这一佐佐木–Einstein 流形及 $AdS_5$ 直积上的 IIB 型超弦理论对偶的四维规范理论，被 Klebanov–Witten [75] 认定为具有超势
\begin{equation*}
\tr\,(A_1B_1A_2B_2-B_1A_1B_2A_2)
\tag{10}
\end{equation*}
的超对称规范理论。

Riemann 锥为环面簇的佐佐木–Einstein 流形称为环面佐佐木–Einstein 流形；在五维时它与 $S^2\times S^3$ 的若干连通和微分同胚。在可能带有环面佐佐木–Einstein 度量的五维流形中，拓扑上仅次于 $S^5$ 之简单的是 $S^2\times S^3$，而 Gauntlett–Martelli–Sparks–Waldram [49] 证明了它带有可数无限多个互不同构的环面佐佐木–Einstein 度量。这一结果是通过把度量具体写成式子来证明的；又，这一结果使人们明白：Cheeger–Tian [29] 曾猜想不存在的、称为非正则（irregular）的那一类佐佐木–Einstein 流形实际上是存在的。这两点值得注目。关于环面佐佐木–Einstein 度量，有
\begin{quote}
\noindent [二小 08] 二木昭人、小野肇，《Einstein 度量与 GIT 稳定性》，数学，Vol.\ 60, No.\ 2, pp.\ 175--202, 2008 年
\end{quote}
可读的日文概说。$(SU(2)\times SU(2))/U(1)$ 与带超势 (10) 的四维超对称规范理论的关系，被推广到一般的环面佐佐木–Einstein 流形；此时对应的规范理论，用称为膜铺砌（brane tiling）的组合论对象来记述。关于这一话题例如见 [144]。这种膜铺砌与 McKay 对应（[107], [117], [119]）有很深的关系。
% [Tutor Note: 校勘。原书「超対称殻ゲージ理論」的「殻」字义不可通（OCR 亦不确定），按上下文（Klebanov–Witten 模型、quiver gauge theory）应读作「超对称规范理论」，或原文残缺一字。中译本按「超对称规范理论」处理。]

% [原文 p.182]
给定三维 Calabi–Yau 流形 $X$ 时，把它变形为具有普通二重点的流形 $X_0$，再取其 crepant 奇点解消，就得到新的 Calabi–Yau 流形 $Y$。$X$ 与 $Y$ 之间的这种关系称为锥流形转迁（conifold transition），在镜像对称中起重要作用，关于此例如见 Batyrev–Ciocan–Fontanine–Kim–van Straten [18]。另一方面，Witten [138] 猜想：$T^*S^3$ 对其零截面的 Lagrange 相交 Floer 理论，给出 $S^3$ 上的 Chern–Simons 理论。如此得到的 $S^3$ 上 Chern–Simons 理论的巨大 $N$ 展开，给出三维普通二重点的 crepant 奇点解消的 Gromov–Witten 不变量——Gopakumar–Vafa [53] 的这一猜想，对此后 Gromov–Witten 理论的发展产生了巨大影响。此外，锥流形转迁对 Reid 的空想（Reid's fantasy）[118]——包含所有三维 Calabi–Yau 流形的既约模空间存在——而言也是本质的。

用 't Hooft [124] 的想法，可以把二维量子引力理论表述为矩阵模型（matrix model）的二重标度极限（double scaling limit）。矩阵模型可以用正交多项式的方法研究，与 KdV 方程的关系可自然地导出。关于这一话题，[141] 有出色的解说。

另一方面，作为二维的某类量子引力理论，Witten 引入了面貌稍异的拓扑引力理论（topological gravity）（[140]）。这一理论的关联函数，由 Riemann 面的模空间（更准确地说，是 Deligne–Mumford 的稳定曲线模空间）上的交点数的母函数给出。进而，Witten 猜想这一拓扑引力理论与先前讨论的量子引力理论等价（[141]）。这一猜想导致惊人的结论：稳定曲线模空间上交点数的母函数是 KdV 方程的 $\tau$ 函数。这就是著名的 Witten 猜想，由 Kontsevich 最先证明（[80]）。Kontsevich 的这一证明，可在
\begin{quote}
\noindent [寺柚 97] 寺柚友秀，《Kontsevich 对 Witten 猜想的解决》，上智大学数学講究録，No.\ 40，1997 年
\end{quote}
读到。此外，弦论家的解说有 [34]。如今这一定理已有种种证明，例如 Okounkov–Pandharipande [114] 的概率论证明、Kazarian–Lando [73] 使用 ELSV 公式 [43] 的证明等。又，Witten 猜想与称为 Virasoro 猜想的另一猜想等价，而 Mirzakhani 用 McShane 恒等式 [92] 的一般化 [96] 直接证明了 Virasoro 猜想（[97]）。

Witten 猜想中出现的稳定曲线模空间的交点数，可以看作目标空间为一点时的 Gromov–Witten 不变量；而 Gromov–Witten 不变量对任意辛流形都有定义，因此对一般的辛流形，其 Gromov–Witten 不变量的母函数也被期待与可积系统有关。特别地，对 $\PP^1$，同变 Gromov–Witten 不变量的母函数成为称为 2-Toda 梯队的可积系统的 $\tau$ 函数，这由 Okounkov–Pandharipande [113] 证明。又，Virasoro 猜想对一般辛流形也被表述（[42]），对环面簇则由 Givental [50] 用二次 Hamilton 函数的量子化（quantization of quadratic Hamiltonian）理论证明。Givental 的这一结果被 Teleman [126] 推广到具有半单量子上同调的任意辛流形。

% [原文 p.183]
此外，对奇点（更准确地说，是复流形与其上全纯函数的组）也期待有 Gromov–Witten 理论的类似物。这原本由 Witten [143] 提出，经众人贡献，对称为可逆多项式的良好一类，已知有 Fan–Jarvis–Ruan [45] 的解析构造与 Polishchuk–Vaintrob [115] 的代数构造。对单纯奇点，这一理论的配分函数给出称为 Drinfeld–Sokolov 梯队（[39]）的可积系统的 $\tau$ 函数。

对 Kähler 流形 $X$ 上的全纯向量丛 $\mathcal{E}$ 与实数 $\mu$，满足
\begin{equation*}
\frac{1}{2\pi\sqrt{-1}}\Lambda F_A=\mu\cdot\id
\tag{11}
\end{equation*}
的联络 $A\in\mathcal{A}^{1,1}$ 称为 $\mathcal{E}$ 的 Hermite Yang–Mills 联络（Hermitian Yang–Mills connection）。当 $X$ 为 Riemann 面时，Hermite Yang–Mills 联络与常中心曲率联络一致。由 Bianchi 恒等式
\begin{equation*}
d_AF_A=\partial_AF_A+\bar\partial_AF_A=0
\tag{12}
\end{equation*}
取 $(2,1)$ 成分与 $(1,2)$ 成分，得
\begin{equation*}
\partial_AF_A=\bar\partial_AF_A=0,
\tag{13}
\end{equation*}
再结合 Kähler 流形上 Hermite 全纯向量丛的恒等式
\begin{align*}
[\Lambda,\partial_A]&=\sqrt{-1}\,\bar\partial_A^*, \tag{14}\\
[\Lambda,\bar\partial_A]&=-\sqrt{-1}\,\partial_A^* \tag{15}
\end{align*}
就有
\begin{align*}
d_A^*F_A&=(\partial_A^*+\bar\partial_A^*)F_A \tag{16}\\
&=\sqrt{-1}\bigl([\Lambda,\bar\partial_A]-[\Lambda,\partial_A]\bigr)F_A \tag{17}\\
&=-\sqrt{-1}(\bar\partial_A-\partial_A)\Lambda F_A, \tag{18}
\end{align*}
故 Hermite Yang–Mills 联络是 Yang–Mills 联络。对 Kähler 流形 $X$ 上的全纯向量丛 $\mathcal{E}$，$\mathcal{E}$ 上的直和不可分解 Hermite Yang–Mills 联络若存在则唯一，且存在的充要条件是 $\mathcal{E}$ 稳定。这是第 9 章主定理 Narasimhan–Seshadri 定理的一般化：对射影曲面由 Donaldson [36] 证明，一般情形由 Uhlenbeck–Yau [131] 证明。这种全纯向量丛与 Hermite Yang–Mills 联络的对应称为小林–Hitchin 对应（Kobayashi–Hitchin correspondence）。
% [Tutor Note: 难点。(11) 中的 $\Lambda$ 是关于 Kähler 形式的内乘（与 $\omega$ 作缩并），$\Lambda F_A=\mu\,\id$ 是「常中心曲率」条件的矩阵形式：$F_A$ 的 $\omega$-迹（即中心部分）为常数 $\mu$ 乘以恒等。$X$ 为曲线时 $F_A$ 是 $(1,1)$ 形式，$\Lambda F_A$ 就是标量曲率，故与「常中心曲率」一致。]

与 Atiyah–Bott [11]、Donaldson [35] 并列，对二维规范理论的发展产生巨大影响的论文，可以举出 Verlinde [134] 与 Witten [142]。Riemann 面上稳定向量丛的模空间，带有称为行列式丛的自然全纯直线丛 $\mathcal{L}$。对自然数 $k$，$\mathcal{L}$ 的 $k$ 阶张量幂 $\mathcal{L}^k$ 的全纯截面空间 $H^0(\mathcal{L}^k)$，与称为 Wess–Zumino–Witten 模型的二维共形场论中 level $k$ 的共形块空间等同。Verlinde（[134]）研究了共形块空间所具有的代数结构，其结果发现了共形块空间维数的公式。这一 Verlinde 公式引起了代数几何学者的注目，已有种种证明。关于 Verlinde 公式的解说文章例如有 [20]。可作为日文阅读的文献，则举出 [向井 08] 与
\begin{quote}
\noindent [上清 08] 上野健尔、清水勇二，《复结构的形变与周期》，岩波书店，2008 年
\end{quote}
Verlinde 公式是直线丛的全纯截面空间的维数公式；由于这一直线丛没有高次上同调，全纯截面空间的维数与 Euler 数一致，依 Riemann–Roch 定理归结为示性类的积分。Witten [142] 给出了模空间上更一般的积分的显式公式。这一 Witten 公式也不亚于 Verlinde 公式而备受注目，已有种种手法的证明，关于此例如见 [5], [70], [94], [127]。

% [原文 p.184]
Riemann 面上秩 2 稳定向量丛的模空间上，存在由 Goldman [52] 构造的完全可积系统。这一完全可积系统与辛环面流形相似，却更为复杂而深奥。这一完全可积系统的动量多面体的格点数，与行列式丛的全纯截面维数一致，并由 Verlinde 公式计算（[71]）。关于这些话题，Audin [14] 的解说易读。此外，作为 Goldman 系的玩具，在多边形空间（polygon space）的模空间上，可以定义称为弯曲系（bending system）的完全可积系统。这与 Guillemin–Sternberg [56] 引入的旗流形上称为 Gelfand–Cetlin 系的完全可积系统、以及后述稳定抛物向量丛模空间上的 Goldman 系，都有很深的关系。

关于把 Riemann 面上稳定向量丛模空间的辛结构在有限维框架中具体记述的问题，已有许多研究；这一方向上的一个决定性工作，当数 Alekseev–Malkin–Meinrenken [6] 的 Lie 群值动量映射（Lie group valued moment map）理论。关于此，有创始人之一 Meinrenken 的讲义录 [95]。

Riemann 面上稳定向量丛的 Narasimhan–Seshadri 定理，被 Mehta–Seshadri [93] 推广到带点 Riemann 面上的稳定抛物向量丛理论。这（几乎）对应于把 Riemann 面推广为轨形（或 Deligne–Mumford 栈），并与用裤子分解（pair-of-pants decomposition）研究 Riemann 面上稳定向量丛的模空间有深刻关联。又，在理解单值保持形变的辛几何侧面时，向抛物向量丛的推广也是本质的。

通过 $2+1=3$ 这一恒等式，二维规范理论与三维规范理论相联结。在三维情形，作为联络空间上的规范不变泛函，除 Yang–Mills 泛函之外还可以考虑 Chern–Simons 泛函
\begin{equation*}
S=\int_X\tr\Bigl(\frac{1}{2}A\wedge dA+\frac{1}{3}A\wedge A\wedge A\Bigr).
\tag{19}
\end{equation*}
这一 Chern–Simons 泛函的驻点由平坦联络组成。把 Chern–Simons 泛函放到指数的肩上、在联络的规范等价类空间上积分所得的拓扑场论研究，由 Witten [139] 开创。三维流形有边界时，对边界确定一个 Hilbert 空间，路径积分给出这一 Hilbert 空间的元；而这一 Hilbert 空间与边界上 Wess–Zumino–Witten 模型的共形块空间等同，其维数由 Verlinde 公式给出。

% [原文 p.185]
Chern–Simons 泛函通过 $4=3+1$ 这一等式，与四维规范理论也密切相关。四维 Riemann 流形上的联络 $A$ 若满足
\begin{equation*}
*F_A=-F_A
\tag{20}
\end{equation*}
则称为反自对偶联络。由 Bianchi 恒等式，反自对偶联络是 Yang–Mills 联络。与 Narasimhan–Seshadri 定理的另一证明同时期，Donaldson 用反自对偶联络的模空间，开发了研究四维流形微分拓扑的全新手法。这就是所谓 Donaldson 理论。作为其结果，具有异种微分结构的四维 Euclid 空间（exotic $\R^4$）的存在性随之推出。

对三维流形 $Y$，Taubes 观察到：$Y$ 上 Chern–Simons 泛函的梯度流轨道，给出四维流形 $Y\times\R$ 上没有 $dt$ 成分的反自对偶联络。这暗示关于 Chern–Simons 泛函的 Morse 理论在带边四维流形的 Donaldson 理论中起重要作用。这样的 Morse 理论由 Floer 定义，被称为 $Y$ 的瞬子 Floer 上同调。关于此见 [深谷 95]。

瞬子 Floer 上同调是对三维流形定义的重要上同调理论；把它与某个 Lagrange 相交 Floer 上同调等同起来，就是 Atiyah（[12]）与 Floer 的著名 Atiyah–Floer 猜想。这一猜想是 Floer 理论中的中心问题之一，朴素地说是指如下之事：由 Morse 理论容易看出，任意紧三维流形 $Y$ 都可以实现为把两个柄体（handlebody）沿边界粘合起来
\begin{equation*}
Y=H_1\cup_\Sigma H_2
\tag{21}
\end{equation*}
这里 $H_1$ 与 $H_2$ 是同胚的柄体，把它们边界上出现的二维流形记为 $\Sigma$，则粘合方式由 $\Sigma$ 的映射类群 $\mathrm{MCG}(\Sigma):=\mathrm{Diff}(\Sigma)/\mathrm{Diff}_0(\Sigma)$ 的元指定。

设 $G$ 为紧 Lie 群，把 $\Sigma$ 上平坦 $G$ 丛的模空间记为 $\mathcal{M}(\Sigma)$。对 $i=1,2$，柄体 $H_i$ 上平坦 $G$ 丛的模空间 $\mathcal{M}(H_i)$ 是 $\mathcal{M}(\Sigma)$ 的 Lagrange 子流形，其公共部分给出 $Y$ 上平坦 $G$ 丛的模空间 $\mathcal{M}(Y)$：
\begin{equation*}
\mathcal{M}(H_1)\cap\mathcal{M}(H_2)=\mathcal{M}(Y).
\tag{22}
\end{equation*}
平坦丛的模空间带有来自可约联络的奇点，不妨宽容地认为其上的 Lagrange 相交 Floer 上同调仍有定义。由 Taubes [125] 的工作，瞬子 Floer 上同调与 Lagrange 相交 Floer 上同调，都在「Euler 数给出称为 Casson 不变量的三维流形不变量」这一意义上，成为 Casson 不变量的范畴化（categorification）。Atiyah–Floer 猜想主张这两个 Floer 上同调同型。部分结果有 Dostoglou–Salamon [38] 与 Muñoz [101], [102]。

% [原文 p.186]
关于紧超 Kähler 流形的概说，有 Huybrechts 的教授资格论文 [66]。另外，Joyce 关于特殊和乐流形的书 [72] 中也讨论了紧超 Kähler 流形。关于非紧全纯辛流形的概说，有 Boalch 的讲义笔记 [23]。

非紧超 Kähler 流形的例子可以举出 ALE 空间。ALE 是渐近局部 Euclid 的（asymptotically locally Euclidean）的缩写。ALE 空间给出 Einstein 方程的解，故有时也称为引力瞬子（gravitational instanton）。

$SL(2,\C)$ 的离散子群称为 Klein 群。特别地，给定 $SL(2,\C)$ 的有限子群 $\Gamma$ 时，$\C^2/\Gamma$ 是带孤立奇点的代数簇，其极小奇点解消是 crepant 的。这样得到的奇点称为 Klein 奇点。特别地，当 $\Gamma$ 为循环群时，Kronheimer 的构造给出 Gibbons–Hawking 与 Hitchin 发现的超 Kähler 流形。第 10 章所介绍的箭图簇，是作为 ALE 空间上 Yang–Mills 瞬子的模空间而被发现的（[83]）。
% [Tutor Note: 术语。原书「箙多様体」即 quiver variety，中译通作「箭图簇」。此处与第 10 章译名统一。]

复流形 $X$ 上全纯向量丛 $\mathcal{E}$ 与 $H^0(\omega_X\otimes\End\mathcal{E})$ 的元 $\Phi$ 的组 $(\mathcal{E},\Phi)$，称为 Higgs 丛（Higgs bundle）。把四维 Yang–Mills 方程作二维的维数约化（dimensional reduction），就得到 Riemann 面上 Higgs 丛 $(\mathcal{E},\Phi)$ 的方程
\begin{equation*}
\begin{cases}
\bar\partial_A\Phi=0,\\
F_A+[\Phi,\Phi^*]=0,
\end{cases}
\tag{23}
\end{equation*}
关于此，Hitchin [62] 已知如下事实：这一方程也有小林–Hitchin 型对应；若组 $(\mathcal{E},\Phi)$ 在适当意义上稳定，则 (23) 的解唯一存在。由 Serre 对偶性
\begin{equation*}
H^0(\omega_X\otimes\End\mathcal{E})\cong\Ext^1(\mathcal{E},\mathcal{E})^*\cong T^*_{[\mathcal{E}]}\mathcal{M},
\tag{24}
\end{equation*}
故若不顾稳定性，Higgs 丛的模空间 $\mathcal{N}$ 可视为全纯向量丛模空间 $\mathcal{M}$ 的余切丛。正如这一点所暗示的，Higgs 丛的模空间带有自然的超 Kähler 结构。Higgs 丛模空间之上，有称为 Hitchin 系（Hitchin system）的完全可积系统结构。Hitchin 系通过特殊化给出种种完全可积系统，不仅与 Seiberg–Witten 理论有深切关联，在 Langlands 对应中也起重要作用（[21], [109]）。超 Kähler 流形带有由 $S^2$ 参数化的复结构族，其中一族给出 Higgs 丛的模空间，其余给出 $X$ 上平坦 $G_\C$ 丛的模空间。
% [Tutor Note: 难点。(24) 中 $\Ext^1(\mathcal{E},\mathcal{E})^*$ 与 $T^*_{[\mathcal{E}]}\mathcal{M}$ 的等同，来自形变理论：$\Ext^1$ 分类 $\mathcal{E}$ 的一阶形变，故是切空间；对偶即余切空间。Higgs 场 $\Phi\in H^0(\omega_X\otimes\End\mathcal{E})$ 由 Serre 对偶正是余切向量，这解释了为何 Higgs 丛模空间 = 余切丛，从而自动带有超 Kähler 结构（余切丛总有 hyperkähler 度量）。]

由 Riemann–Hilbert 对应，平坦 $G_\C$ 丛的模空间与全纯联络的模空间等同。全纯联络模空间上的全纯辛结构，正是以 Painlevé 方程为代表的单值保持形变所具有的 Hamilton 结构的由来。关于单值保持形变的辛几何侧面，见 Hitchin 的解说 [64]。

单值保持形变通过 Frobenius 流形与镜像对称深切相关。关于 Frobenius 流形，有 Dubrovin 的讲义录 [40], [41] 与 Manin 的入门书 [87]。
% [原文 p.187]
镜像对称的入门书以 Cox–Katz [31] 闻名，此外还有 2000 年在 Boston 郊外举行的春季学校的记录 [65]，以及
\begin{quote}
\noindent [秦泉寺 14] 秦泉寺雅夫，《为数学物理学生的镜像对称入门》，サイエンス社（临时别册·数理科学 SGC-109），2014 年
\end{quote}
等。
